\section{任务、随机化与评估细节}
\begin{strip}
\centering
\begin{table}[H]
\centering
\caption{真机任务及语言指令。}
\label{tab:task-instructions}
\small
\setlength{\tabcolsep}{4pt}
\renewcommand{\arraystretch}{1.2}
\begin{tabularx}{\linewidth}{@{}lY@{}}
\toprule
任务 & 指令 \\
\midrule
放入抽屉 & 拿起物体并将其放入抽屉。 \\
放进盒子 & 拿起物体并将其放入盒子。 \\
积木分类 & 将所有积木放入颜色对应的容器。 \\
插接 & 将插接件插入对应孔位。 \\
\bottomrule
\end{tabularx}
\end{table}

\end{strip}
\subsection{示范与训练设置}
每个任务的示范量为\pending；记录 episode 列表、训练分区、相机视角与帧率、动作表示和归一化方式。所有方法共享示范清单；各方法的预训练权重、微调范围、训练步数、batch size 和随机种子填入复现记录。

\subsection{随机化支持范围}
\textbf{位置：}训练和评估均在指定桌面区域内随机，范围、朝向和采样方式为\pending。\textbf{桌布：}训练桌布集合为\pending，OOD 使用的留出桌布编号为\pending。\textbf{物体外观：}目标物体的身份、颜色／纹理版本和保持不变的语义为\pending。\textbf{干扰物：}训练场景配置为\pending，测试留出的干扰物种类、数量和区域为\pending。

每次 OOD 试验只改变一个因素，并复用相同的初始化编号供三个方法比较。道具变化以照片和 ID 记录。

\subsection{试验数、时限与汇总}
当前约定每任务、每方法 20 次，单次 120 秒；四条件分配为\pending。若每条件 20 次，四任务、三方法、四条件共 960 次；若四条件合计 20 次，共 240 次。最终使用的口径在采集清单和主表中保持一致。

成功率 $\mathrm{SR}=N_{\mathrm{success}}/N_{\mathrm{trial}}$。每次试验使用表~\ref{tab:criteria} 的完整终点判据。保留原始 $k/n$，并按任务等权汇总；阶段完成情况作为诊断字段单独记录。

\subsection{结束条件与失败类型}
达到完整终点时记录成功和完成时间；到达 120 秒时记录超时。人工中止、设备异常等事件分别登记，并在汇总前依据统一规则处理。失败类型包括抓取失败、放置越界、对齐失败、碰撞、目标混淆与超时，按实际试验补充。

\subsection{本附录对应的参考}
Efficient-WAM 在正文 \S4.3 写示范量、试验数、时限和五任务判据；WAM4D Table 4 列任务分步判据；OpenWAM Appendix C 提供硬件图、任务序列和终点定义。本项目以一张任务表加简短协议文字组织这些材料。
